% Shared activity layout. Loaded by utpnotes.tex.
% A printed output and its "Reference output." label stay together on one page.
\newenvironment{refoutput}{\par\noindent\begin{minipage}{\linewidth}\expectedoutput}{\end{minipage}\par}

% Start a new page unless at least #1 of room is left, so a box never opens with one stranded line.
\newcommand{\keepspace}[1]{\par\penalty-100\relax\ifdim\dimexpr\pagegoal-\pagetotal\relax<#1\newpage\fi}

% Activity boxes shared by notebook-based labs.
\utpnewcallout[boxrule=0.9pt]{yourturn}{white}{UTPTeal}{Your turn}
\utpnewcallout[boxrule=0.65pt,left=2.5mm,right=2.5mm,top=1.6mm,bottom=1.6mm,colbacktitle=UTPMuted]{expectedfindings}{UTPPaleGray!72}{UTPMuted!62}{Expected findings}

% A band that opens a part of the notebook inside a lab sheet.
\newcommand{\labpart}[1]{\keepspace{11\baselineskip}\par\medskip\noindent{\large\bfseries\color{UTPBlue}#1}\par\nopagebreak\smallskip}

% ------------------------------------------------------------- lab parts
% A lab is worked in parts. The map comes before the first question; each part opens with a
% banner giving its goal, questions and data (on a new page unless most of the page is free),
% and closes with a checkpoint.
%
%   \labroadmap{This lab in four parts}{A & What you do & 2--5 \\ ... }
%   \labpartopen{Part A}{Title}{Questions 2--5 \textbullet\ Data: ...}{One-sentence goal.}
%   \begin{labcheckpoint} \begin{itemize} \item ... \end{itemize} \end{labcheckpoint}
%   \labsubhead{Try it yourself}
\newcommand{\labroadmap}[2]{%
  \par\smallskip\noindent
  \begin{tcolorbox}[utpplain,unbreakable,colback=white,colframe=UTPNavy,boxrule=.7pt,arc=2mm,
      left=2.5mm,right=2.5mm,top=1.4mm,bottom=1.4mm,fonttitle=\bfseries\color{white},colbacktitle=UTPNavy,title={#1}]
    \small\renewcommand{\arraystretch}{1.22}%
    \begin{tabularx}{\linewidth}{@{}>{\bfseries\leavevmode\color{UTPNavy}}p{1.35cm}X>{\raggedleft\arraybackslash}p{1.55cm}@{}}
      {\color{UTPMuted}Part} & \leavevmode{\bfseries\color{UTPMuted}What you do} & \leavevmode{\bfseries\color{UTPMuted}Questions}\\[1pt]
      \hline
      #2
    \end{tabularx}
  \end{tcolorbox}\par\smallskip}

\newcounter{labpartno}
\newcommand{\labpartopen}[4]{%
  \par\bigskip
  \keepspace{0.6\textheight}% a part opens a page, unless most of this one is still free
  \stepcounter{labpartno}%
  \pdfbookmark[3]{#1: #2}{labpart.\arabic{labpartno}}%
  \noindent
  \begin{tcolorbox}[utpplain,unbreakable,colback=UTPPaleBlue,colframe=UTPNavy,boxrule=.8pt,arc=2mm,
      left=3mm,right=3mm,top=2mm,bottom=2mm,before skip=0pt,after skip=1em,
      fonttitle=\large\bfseries\color{white},colbacktitle=UTPNavy,toptitle=1.6mm,bottomtitle=1.6mm,
      title={#1\hspace{.55em}{\color{UTPTeal!45!white}\textbar}\hspace{.55em}#2}]
    \textbf{\color{UTPNavy}Goal.} #4\par\smallskip
    {\small\color{UTPMuted}#3\par}
  \end{tcolorbox}\par}

\utpnewcallout[unbreakable,boxrule=0.8pt,before skip=1.2em]{labcheckpoint}{UTPPaleTeal!65}{UTPGreen}{Checkpoint: before you move on}

\newcommand{\labsubhead}[1]{\keepspace{9\baselineskip}\par\medskip\noindent{\large\bfseries\color{UTPBlue}#1}\par\nopagebreak\smallskip}

% The verification and submission checklists stay together on one page.
\newenvironment{labchecklists}{\par\noindent\begin{minipage}{\linewidth}\tcbset{break at=100cm}}{\end{minipage}\par}


% Shared notebook routine; courses may override its introductory and comparison text.
\providecommand{\utplabroutineintro}{For each question, do four things in order:}
\providecommand{\utplabroutinecompare}{Check your output against the reference output.}
\newcommand{\labroutine}{%
  \par\smallskip\noindent\textbf{How to work through this sheet.} \utplabroutineintro
  \begin{utpsteps}
    \item \textbf{Predict.} Read the task and write down what you expect before running anything.
    \item \textbf{Run.} Run the matching cell in the notebook. The listing printed under the question is that same code, so you can read it here and run it there.
    \item \textbf{Compare.} \utplabroutinecompare
    \item \textbf{Explain.} Add a Markdown cell saying what the result means, in your own words. Where your first prediction was wrong, say why.
  \end{utpsteps}
  \par\smallskip}
